\documentclass[12pt,a4paper]{report}

% ---------- Paquets ----------
\usepackage[utf8]{inputenc}
\usepackage[T1]{fontenc}
\usepackage[french]{babel}
\usepackage{lmodern}
\usepackage[margin=2.5cm]{geometry}
\usepackage{amsmath}
\usepackage{graphicx}
\usepackage{booktabs}
\usepackage{tabularx}
\usepackage{xcolor}
\usepackage{tikz}
\usetikzlibrary{arrows.meta,positioning,shapes,shapes.multipart,babel}
\pgfdeclarelayer{background}
\pgfsetlayers{background,main}
\usepackage{fancyhdr}
\usepackage{enumitem}
\usepackage[hidelinks]{hyperref}

% ---------- En-têtes et pieds de page ----------
\pagestyle{fancy}
\fancyhf{}
\fancyhead[L]{Cahier des Charges}
\fancyhead[R]{Sawti-ASR-FineTuner \& Data-Augmenter}
\fancyfoot[C]{\thepage}
\renewcommand{\headrulewidth}{0.4pt}
\setlength{\headheight}{15pt}

\begin{document}

% ---------- Page de garde ----------
\begin{titlepage}
  \centering
  {\large École Nationale Supérieure d'Arts et Métiers -- Meknès\par}
  {\large Filière Intelligence Artificielle -- 4\textsuperscript{e} année\par}
  \vspace{3cm}
  {\Huge\bfseries Cahier des Charges\par}
  \vspace{1cm}
  {\LARGE Sawti-ASR-FineTuner \& Data-Augmenter\par}
  \vspace{0.5cm}
  {\large Pipeline ASR \& MLOps Audio pour la Darija marocaine\\
  et le code-switching français--arabe\par}
  \vfill
  \begin{tabular}{ll}
    \textbf{Réalisé par :} & Mohamed Amine Belasri \\
                           & Yahya Amajane \\[0.5cm]
    \textbf{Encadré par :} & Pr. Hosni \\
                           & Pr. Menaoui \\
  \end{tabular}
  \vfill
  {\large Année universitaire 2026--2027\par}
\end{titlepage}

\tableofcontents
\listoftables
\listoffigures

% ---------- Parties (à rédiger une par une) ----------
% =====================================================================
\chapter{Présentation du projet}
% =====================================================================

\section{Titre du projet}
\textbf{Sawti-ASR-FineTuner \& Data-Augmenter} -- Pipeline ASR \& MLOps Audio.

Il s'agit d'un projet d'initiation de 4\textsuperscript{e} année (filière IA, ENSAM Meknès),
réalisé en binôme sur un semestre (16 semaines), dans le cadre de la plateforme
Sawti (thème : continuité, Arabizi et code-switching).

\section{Problème à résoudre}
Les modèles fondationnels de reconnaissance automatique de la parole (ASR), comme
Whisper, sont entraînés principalement sur des langues bien dotées en données. Ils
donnent de mauvais résultats sur la Darija marocaine, pour trois raisons :
\begin{itemize}
  \item \textbf{Rareté des données} : il existe peu de corpus de Darija transcrits.
  \item \textbf{Code-switching} : les locuteurs alternent la Darija et le français
        dans une même phrase.
  \item \textbf{Conditions réelles} : les enregistrements contiennent du bruit de fond
        (rue, café, téléphone), absent des corpus propres.
\end{itemize}
De plus, adapter ces modèles par un \emph{full fine-tuning} coûte cher en calcul,
alors que le projet ne dispose que de GPU gratuits (Kaggle et Google Colab, T4/P100).

\section{Objectif de la solution}
Le projet vise à concevoir un pipeline MLOps automatisé qui :
\begin{enumerate}
  \item augmente les données audio (SpecAugment, injection de bruit, variations de
        pitch et de vitesse) pour compenser leur rareté ;
  \item fine-tune efficacement le modèle Whisper avec LoRA (PEFT) ;
  \item évalue la robustesse du modèle sur trois jeux de test : propre, bruité et
        code-switching ;
  \item trace toutes les expériences et versionne les modèles avec MLflow.
\end{enumerate}
L'objectif mesurable est de battre la baseline (full fine-tuning sans augmentation
ni PEFT) en gain relatif de WER. Les valeurs cibles sont données dans la partie~5.

\section{Périmètre}
Le tableau~\ref{tab:perimetre} précise ce que le projet fait et ne fait pas.

\begin{table}[h]
  \centering
  \small
  \begin{tabularx}{\textwidth}{XX}
    \toprule
    \textbf{Inclus dans le projet} & \textbf{Exclu du projet} \\
    \midrule
    Préparation des données audio au format Hugging Face Datasets (nettoyage,
    rééchantillonnage à 16~kHz, découpage train/validation/test).
    & Collecte ou enregistrement de nouveaux audios sur le terrain. \\
    \addlinespace
    Pipeline d'augmentation acoustique et spectrale avec niveaux de bruit
    contrôlés (SNR).
    & Transcription manuelle à grande échelle de nouvelles données. \\
    \addlinespace
    Baseline : Whisper zero-shot et full fine-tuning (version réduite compatible
    avec un GPU T4).
    & Interface utilisateur ou application finale de transcription. \\
    \addlinespace
    Fine-tuning LoRA/PEFT et étude de ses hyperparamètres (rang, modules ciblés).
    & Déploiement en production (serveur, API temps réel). \\
    \addlinespace
    Évaluation WER/CER sur trois jeux de test et calcul du gain relatif.
    & Entraînement de grands modèles (Whisper medium/large) ou depuis zéro. \\
    \addlinespace
    Suivi des expériences, registre de modèles (MLflow) et script de
    réentraînement automatisé.
    & Synthèse vocale (TTS) et traduction. \\
    \bottomrule
  \end{tabularx}
  \caption{Matrice d'inclusion et d'exclusion du projet}
  \label{tab:perimetre}
\end{table}

\section{Utilisateurs visés}
\begin{itemize}
  \item \textbf{Équipe de la plateforme Sawti et encadrants} : consultent les
        résultats et comparent les modèles.
  \item \textbf{Ingénieurs ML / data scientists} : lancent le pipeline, ajoutent
        des données et réentraînent les modèles.
  \item \textbf{Communauté open source} (utilisateurs secondaires) : réutilisent le
        code publié sur GitHub et le modèle publié sur Hugging Face.
\end{itemize}

% =====================================================================
\chapter{Exigences fonctionnelles}
% =====================================================================

\section{Fonctionnalités principales}
\begin{table}[h]
  \centering
  \small
  \begin{tabularx}{\textwidth}{lp{3.6cm}X}
    \toprule
    \textbf{ID} & \textbf{Fonctionnalité} & \textbf{Description} \\
    \midrule
    F1 & Préparation des données & Charger les corpus, nettoyer, rééchantillonner à
    16~kHz, normaliser les transcriptions, découper en train/validation/test. \\
    F2 & Augmentation & Ajouter du bruit (ESC-50, NOISEX) à un SNR contrôlé, faire
    varier le pitch et la vitesse, appliquer SpecAugment. \\
    F3 & Baseline & Évaluer Whisper en zero-shot et réaliser un full fine-tuning
    sans augmentation. \\
    F4 & Fine-tuning LoRA & Entraîner des adaptateurs LoRA sur Whisper, avec ou
    sans augmentation. \\
    F5 & Évaluation & Calculer le WER et le CER sur les jeux propre, bruité et
    code-switching, puis le gain relatif par rapport à la baseline. \\
    F6 & Suivi des expériences & Enregistrer dans MLflow les paramètres, les
    métriques et les artefacts de chaque expérience. \\
    F7 & Registre de modèles & Versionner les modèles et identifier la meilleure
    version. \\
    F8 & Réentraînement automatisé & Relancer toute la chaîne avec une seule
    commande et un fichier de configuration. \\
    F9 & Publication & Publier le meilleur modèle et sa documentation sur
    Hugging Face. \\
    \bottomrule
  \end{tabularx}
  \caption{Fonctionnalités principales}
  \label{tab:fonctionnalites}
\end{table}

\section{Entrées attendues}
\begin{itemize}
  \item Fichiers audio (WAV ou FLAC, mono, 16~kHz) accompagnés de leur transcription.
  \item Corpus de bruits (ESC-50, NOISEX) pour l'augmentation.
  \item Fichier de configuration YAML : modèle, hyperparamètres, niveaux de SNR.
  \item Checkpoints existants, pour reprendre un entraînement interrompu.
\end{itemize}

\section{Sorties attendues}
\begin{itemize}
  \item Modèles fine-tunés et adaptateurs LoRA (checkpoints).
  \item Transcriptions prédites sur les jeux de test.
  \item Métriques WER et CER par jeu de test, et tableau comparatif des expériences.
  \item Runs MLflow (paramètres, courbes, artefacts) et versions de modèles dans
        le registre.
\end{itemize}

\section{Scénarios d'utilisation}
\paragraph{Scénario 1 -- Expérience complète.}
L'ingénieur rédige un fichier de configuration et lance le pipeline. Les données
sont préparées puis augmentées, le modèle est entraîné avec des checkpoints
réguliers, puis évalué sur les trois jeux de test. Les résultats sont enregistrés
dans MLflow. Si le modèle est meilleur que la version actuelle, il est ajouté au
registre.

\paragraph{Scénario 2 -- Réentraînement avec de nouvelles données.}
De nouvelles données sont ajoutées au corpus. L'ingénieur relance la même commande.
Une nouvelle version du modèle est produite et comparée à l'ancienne.

\paragraph{Scénario 3 -- Reprise après une coupure du GPU.}
La session Kaggle ou Colab s'arrête pendant l'entraînement. L'ingénieur relance le
pipeline, qui reprend au dernier checkpoint sauvegardé.

\paragraph{Scénario 4 -- Comparaison des expériences.}
L'encadrant ouvre MLflow et compare les runs (ablations : avec ou sans augmentation,
LoRA ou full fine-tuning, différents rangs et niveaux de SNR).

\section{Priorité des fonctionnalités}
\begin{table}[h]
  \centering
  \small
  \begin{tabular}{lll}
    \toprule
    \textbf{ID} & \textbf{Fonctionnalité} & \textbf{Priorité} \\
    \midrule
    F1 & Préparation des données    & Indispensable \\
    F2 & Augmentation               & Indispensable \\
    F3 & Baseline                   & Indispensable \\
    F4 & Fine-tuning LoRA           & Indispensable \\
    F5 & Évaluation                 & Indispensable \\
    F6 & Suivi des expériences      & Importante \\
    F7 & Registre de modèles        & Importante \\
    F8 & Réentraînement automatisé  & Importante \\
    F9 & Publication                & Optionnelle \\
    \bottomrule
  \end{tabular}
  \caption{Priorité des fonctionnalités}
  \label{tab:priorites}
\end{table}

% =====================================================================
\chapter{Diagrammes UML}
% =====================================================================

\section{Diagramme de cas d'utilisation}
L'ingénieur ML utilise toutes les fonctions du pipeline. L'encadrant consulte et
compare les résultats. MLflow est un acteur système qui stocke les expériences et
les versions de modèles.

\begin{figure}[h]
  \centering
  \begin{tikzpicture}[
      uc/.style={draw, ellipse, fill=white, minimum width=4.2cm, minimum height=0.85cm,
                 font=\small, align=center},
      actor/.style={draw, rectangle, rounded corners, fill=gray!15,
                    minimum width=2.6cm, minimum height=0.9cm,
                    font=\small, align=center}]
    \draw[thick] (-0.5,-0.2) rectangle (6.5,-9.8);
    \node at (3,-0.55) {\textbf{Pipeline Sawti-ASR}};
    \node[uc] (u1) at (3,-1.5) {Préparer les données};
    \node[uc] (u2) at (3,-2.6) {Augmenter les données};
    \node[uc] (u3) at (3,-3.7) {Entraîner la baseline};
    \node[uc] (u4) at (3,-4.8) {Fine-tuner avec LoRA};
    \node[uc] (u5) at (3,-5.9) {Évaluer (WER / CER)};
    \node[uc] (u6) at (3,-7.0) {Comparer les expériences};
    \node[uc] (u7) at (3,-8.1) {Gérer le registre};
    \node[uc] (u8) at (3,-9.2) {Relancer le réentraînement};
    \node[actor] (a1) at (-3.1,-5.0) {Ingénieur ML};
    \node[actor] (a2) at (9.1,-7.5) {Encadrant};
    \node[actor] (a3) at (9.1,-3.5) {MLflow\\(système)};
    \begin{pgfonlayer}{background}
      \foreach \u in {u1,u2,u3,u4,u5,u6,u7,u8} \draw (a1) -- (\u);
      \draw (a2) -- (u6);
      \draw (a2) -- (u7);
      \draw (a3) -- (u5);
      \draw (a3) -- (u6);
      \draw (a3) -- (u7);
    \end{pgfonlayer}
  \end{tikzpicture}
  \caption{Diagramme de cas d'utilisation}
  \label{fig:cas-utilisation}
\end{figure}

\clearpage
\section{Diagramme d'activité}
Le diagramme suivant décrit le déroulement d'une exécution complète du pipeline.

\begin{figure}[h]
  \centering
  \begin{tikzpicture}[
      node distance=0.5cm, >=Stealth,
      act/.style={draw, rounded corners=6pt, minimum width=6.5cm,
                  minimum height=0.8cm, font=\small, align=center},
      dec/.style={draw, diamond, aspect=2.5, font=\small, align=center,
                  inner sep=1pt}]
    \node[circle, fill=black, minimum size=0.4cm] (s) {};
    \node[act, below=of s]  (a1) {Charger la configuration (YAML)};
    \node[act, below=of a1] (a2) {Charger et nettoyer les données\\(16 kHz, normalisation du texte)};
    \node[act, below=of a2] (a3) {Découper en train / validation / test};
    \node[act, below=of a3] (a4) {Augmenter les données d'entraînement\\(bruit, pitch, vitesse, SpecAugment)};
    \node[act, below=of a4] (a5) {Entraîner le modèle (baseline ou LoRA)\\avec checkpoints réguliers};
    \node[act, below=of a5] (a6) {Évaluer sur les trois jeux de test};
    \node[act, below=of a6] (a7) {Enregistrer le run dans MLflow};
    \node[dec, below=of a7] (d)  {Meilleur que la\\version actuelle ?};
    \node[act, below=0.8cm of d] (a8) {Enregistrer une nouvelle version\\dans le registre};
    \node[circle, draw, minimum size=0.55cm, below=of a8] (e) {};
    \node[circle, fill=black, minimum size=0.3cm] at (e) {};
    \draw[->] (s)  -- (a1);
    \draw[->] (a1) -- (a2);
    \draw[->] (a2) -- (a3);
    \draw[->] (a3) -- (a4);
    \draw[->] (a4) -- (a5);
    \draw[->] (a5) -- (a6);
    \draw[->] (a6) -- (a7);
    \draw[->] (a7) -- (d);
    \draw[->] (d)  -- node[right] {oui} (a8);
    \draw[->] (a8) -- (e);
    \draw[->] (d.east) -- node[above] {non} ++(2,0) |- (e.east);
  \end{tikzpicture}
  \caption{Diagramme d'activité du pipeline}
  \label{fig:activite}
\end{figure}

\clearpage
\section{Diagramme de séquence}
Le scénario principal est l'expérience complète (scénario~1).

\begin{figure}[h]
  \centering
  \begin{tikzpicture}[
      >=Stealth,
      obj/.style={draw, fill=gray!15, minimum width=2.6cm, minimum height=0.8cm,
                  font=\small, align=center},
      lbl/.style={font=\scriptsize, fill=white, inner sep=1pt}]
    \node[obj] (i) at (0,0)    {Ingénieur ML};
    \node[obj] (p) at (3.2,0)  {Pipeline};
    \node[obj] (t) at (6.4,0)  {Entraîneur\\(Whisper + LoRA)};
    \node[obj] (v) at (9.6,0)  {Évaluateur};
    \node[obj] (m) at (12.8,0) {MLflow};
    \foreach \x in {0,3.2,6.4,9.6,12.8} \draw[dashed] (\x,-0.6) -- (\x,-9.6);
    \draw[->] (0,-1.2) -- node[lbl, above] {lancer(config)} (3.2,-1.2);
    \draw[->] (3.2,-1.8) -- (3.9,-1.8) -- (3.9,-2.5) -- (3.2,-2.5);
    \node[font=\scriptsize, align=left, anchor=west] at (3.95,-2.15) {préparer\\augmenter};
    \draw[->] (3.2,-3.1) -- node[lbl, above] {entraîner(données, config)} (6.4,-3.1);
    \draw[->, dashed] (6.4,-3.9) -- node[lbl, above] {modèle entraîné} (3.2,-3.9);
    \draw[->] (3.2,-4.7) -- node[lbl, above, pos=0.7] {évaluer(modèle, 3 jeux de test)} (9.6,-4.7);
    \draw[->, dashed] (9.6,-5.5) -- node[lbl, above, pos=0.3] {WER / CER} (3.2,-5.5);
    \draw[->] (3.2,-6.3) -- node[lbl, above, pos=0.75] {enregistrer paramètres et métriques} (12.8,-6.3);
    \draw[->] (3.2,-7.1) -- node[lbl, above, pos=0.75] {enregistrer la version (si meilleure)} (12.8,-7.1);
    \draw[->, dashed] (12.8,-7.9) -- node[lbl, above, pos=0.25] {numéro de version} (3.2,-7.9);
    \draw[->, dashed] (3.2,-8.8) -- node[lbl, above] {rapport} (0,-8.8);
  \end{tikzpicture}
  \caption{Diagramme de séquence de l'expérience complète}
  \label{fig:sequence}
\end{figure}

\section{Diagramme d'états-transitions}
L'entité qui change d'état est la \textbf{version de modèle} : elle est entraînée,
évaluée, puis rejetée ou promue, et enfin archivée lorsqu'une meilleure version la
remplace. Dans MLflow, ces états seront représentés par des étiquettes ou des alias
de version.

\begin{figure}[h]
  \centering
  \begin{tikzpicture}[
      >=Stealth,
      st/.style={draw, rounded corners=8pt, minimum width=2.4cm,
                 minimum height=0.8cm, font=\small},
      lbl/.style={font=\scriptsize, align=center}]
    \node[circle, fill=black, minimum size=0.35cm] (s) at (0,0) {};
    \node[st] (e1) at (2.4,0)   {Entraîné};
    \node[st] (e2) at (6.2,0)   {Évalué};
    \node[st] (e3) at (10.4,0)  {Candidat};
    \node[st] (e4) at (10.4,-2.6) {Production};
    \node[st] (e5) at (6.2,-2.6)  {Archivé};
    \node[st] (e6) at (6.2,2.4)   {Rejeté};
    \node[circle, draw, minimum size=0.5cm] (f1) at (2.4,-2.6) {};
    \node[circle, fill=black, minimum size=0.28cm] at (f1) {};
    \node[circle, draw, minimum size=0.5cm] (f2) at (9.4,2.4) {};
    \node[circle, fill=black, minimum size=0.28cm] at (f2) {};
    \draw[->] (s)  -- (e1);
    \draw[->] (e1) -- node[lbl, above] {évaluation\\terminée} (e2);
    \draw[->] (e2) -- node[lbl, above] {WER $<$\\version actuelle} (e3);
    \draw[->] (e2) -- node[lbl, left] {WER $\geq$ version actuelle} (e6);
    \draw[->] (e6) -- (f2);
    \draw[->] (e3) -- node[lbl, right] {validation} (e4);
    \draw[->] (e4) -- node[lbl, above] {nouvelle version\\meilleure} (e5);
    \draw[->] (e5) -- (f1);
  \end{tikzpicture}
  \caption{Diagramme d'états-transitions d'une version de modèle}
  \label{fig:etats}
\end{figure}

\clearpage
\section{Diagramme de classes conceptuel}
Le diagramme présente les entités métier du projet et leurs relations.

\begin{figure}[h]
  \centering
  \begin{tikzpicture}[
      cls/.style={draw, rectangle split, rectangle split parts=2,
                  font=\footnotesize, align=left, minimum width=3.4cm},
      mul/.style={font=\scriptsize}]
    \node[cls] (corpus) at (0,0) {\textbf{Corpus}
      \nodepart{second} nom\\source\\volume (heures)};
    \node[cls] (enr) at (5.2,0) {\textbf{EnregistrementAudio}
      \nodepart{second} identifiant\\durée (s)\\fréquence (16 kHz)\\split};
    \node[cls] (tr) at (10.4,0) {\textbf{Transcription}
      \nodepart{second} texte\\code-switching (oui/non)};
    \node[cls] (exp) at (0,-4.2) {\textbf{Expérience}
      \nodepart{second} identifiant du run\\type (baseline / LoRA)\\hyperparamètres\\date};
    \node[cls] (aug) at (5.2,-4.2) {\textbf{Augmentation}
      \nodepart{second} type\\SNR (dB)\\paramètres};
    \node[cls] (mod) at (0,-8.2) {\textbf{Modèle}
      \nodepart{second} modèle de base (Whisper)\\taille\\rang LoRA};
    \node[cls] (ver) at (5.2,-8.2) {\textbf{VersionModèle}
      \nodepart{second} numéro\\état};
    \node[cls] (met) at (10.4,-7.6) {\textbf{Métrique}
      \nodepart{second} jeu de test\\WER\\CER};
    \draw (corpus) -- node[mul, above, pos=0.15] {1} node[mul, above, pos=0.85] {*} (enr);
    \draw (enr) -- node[mul, above, pos=0.15] {1} node[mul, above, pos=0.85] {1} (tr);
    \draw (aug) -- node[mul, right, pos=0.15] {*} node[mul, right, pos=0.85] {*} (enr);
    \draw (exp) -- node[mul, right, pos=0.15] {*} node[mul, right, pos=0.85] {1} (corpus);
    \draw (exp) -- node[mul, above, pos=0.15] {*} node[mul, above, pos=0.85] {*} (aug);
    \draw (exp) -- node[mul, right, pos=0.15] {1} node[mul, right, pos=0.85] {1} (mod);
    \draw (mod) -- node[mul, above, pos=0.15] {1} node[mul, above, pos=0.85] {*} (ver);
    \draw (exp) -- node[mul, above, pos=0.1] {1} node[mul, above, pos=0.9] {*} (met);
  \end{tikzpicture}
  \caption{Diagramme de classes conceptuel}
  \label{fig:classes}
\end{figure}

\chapter{Données}
\textit{(À rédiger.)}

\chapter{Approche IA}
\textit{(À rédiger.)}

\chapter{Exigences non fonctionnelles}
\textit{(À rédiger.)}

\chapter{Planning et livrables}
\textit{(À rédiger.)}

\chapter{Annexes}
\textit{(À rédiger.)}

\end{document}
